\begin{abstract}
Text-to-image models fail at the smallest images people draw on purpose: 12--24\,px RGBA game sprites, where a
handful of colours and hard alpha edges \emph{are} the medium. We show that the usual failure is not a modelling
failure but a target failure. In the standard sprite corpus only 7\% of the images are natively 16\,px or smaller;
the rest is 25--64\,px artwork box-downscaled into the low-resolution buckets, so a model trained on it learns soft,
fifty-colour blobs---and the evaluation reference inherits the same bias, rewarding exactly that. We rebuild the
targets (drop 4{,}888 pixel-identical duplicates that leak 15.5\% of the held-out set into training, forbid
downscaling by more than $1.5\times$, add 14k natively small sprites from three permissively licensed corpora) and we
evaluate against a reference drawn only from real low-resolution pixel art. Under this protocol the data fix alone
takes FD-DINOv2 from 209.6 to 117.9 at 16\,px. We then close the remaining gap where it actually is: real sprites use
a median of six opaque colours while every sampling configuration we tested emits more than fifty. Projecting the
predicted $x_0$ onto a per-sprite $k$-colour palette during the late denoising steps---so that the remaining steps
repair the quantisation seams---reaches $31.8\pm0.5$, better than quantising the finished samples (41.6), and
$2.6\times$ better than the strongest external system under the same palette post-process (50.8 versus 134.4 at
$n{=}200$); it also beats the mixed-provenance ``real'' sprites that the downscale-biased protocol treats as ground
truth (77.0). We report an equally instructive negative result: cross-resolution self-guidance and its internalised
one-pass variant, which win by 40\% under the downscale-biased reference, give no gain once the reference is real
pixel art, while independent-condition guidance transfers and composes with palette projection.
\end{abstract}
